% --- ТОВЧЛОЛ ---
\newpage
\chapter*{Товчлол}
\addcontentsline{toc}{chapter}{Товчлол}

Өвөрхангай аймгийн Хужирт сумын сувиллуудын захиалгын үйл явцыг цахимжуулах зорилгоор веб суурьтай бүртгэл мэдээллийн системийг хөгжүүлж, туршилтын орчинд нэвтрүүлсний үр дүнг энэ судалгааны ажилд тусгав.

\textbf{Гол дэвшүүлсэн санаа:} Сувиллын утсаар захиалга авах, цаасан бүртгэл хөтлөх уламжлалт аргыг орлуулах цахим платформыг Next.js, Prisma ORM, PostgreSQL технологийн стек дээр суурилан хөгжүүлсэн. Систем нь олон хэрэглэгчийн үүрэг (зочин, үйлчлүүлэгч, ажилтан, админ)-тэй, аюулгүй байдал, зэрэг хандалтын хяналттай бүрэн хэмжээний веб системийг хэрэгжүүлсэн.

\textbf{Судалгааны үр дүн:}

Нэгдүгээрт, гүйцэтгэлийн хэмжилтээс харахад шинэ систем нэг захиалга боловсруулах хугацааг 15 минутаас 2.3 минут болгон буюу \textbf{85\%}-иар бууруулсан. Өдрийн тайлан гаргах үйлдэл 35 минутаас 1.2 минут болж \textbf{96.6\%}-ийн бууралттай байлаа. Энэ нь сувиллын ажилтнуудын өдөр тутмын ажлын ачааллыг мэдэгдэхүйц хөнгөвчилдгийг нотолж байна.

Хоёрдугаарт, цаасан бүртгэлд тогтмол тохиолддог байсан давхар захиалгын асуудлыг PostgreSQL-ийн pessimistic locking (\texttt{SELECT ... FOR UPDATE}) болон EXCLUDE constraint-ийн хослолоор бүрэн шийдсэн. 50 хэрэглэгчийн зэрэг хандалтын туршилтад \textbf{нэг ч давхар захиалга} үүсгэлгүй систем зөв ажилласан.

Гуравдугаарт, $n=13$ хэрэглэгчид System Usability Scale \cite{brooke1996sus} аргачлалаар хийсэн үнэлгээгээр системийн ашиглах чадвар \textbf{76.3 оноо} авч, Bangor нарын \cite{bangor2008sus} ``сайн'' ангиллын (68+) босгыг давлаа. Ажилтнуудын дундаж оноо 81.7 байсан нь энгийн, ойлгомжтой интерфейстэй болохыг илтгэж байна.

\textbf{Ач холбогдол:} Энэхүү судалгааны ажлын үр дүн нь Хужирт сумын сувиллуудад практикт нэвтрүүлэх бэлэн жишиг загвар болохоос гадна ижил төрлийн аялал жуулчлал, рашаан сувилалын байгууллагуудад хэвшүүлэх боломжтой. Системийг цаашид гар утасны аппликейшн, банкны QPay/SocialPay интеграци, олон сувиллын (multi-tenant) архитектур руу өргөтгөх боломжтой гэдгийг судлаач тэмдэглэж байна.
